\section{Methodology} \label{sec:methodology}

\subsection{From Constructive Evidence to Temporal Assertions}
Figure~\ref{fig:comparison} separates the source of functional intent from the
reference used for implementation checking. The IOR $R$ generates
$\mathrm{FSR}(R)$ and its characteristic fragment $\Phi_R$. A chosen
$D_{\mathrm{ref}}\in\mathrm{FSR}(R)$ supplies the system $Z_{\mathrm{ref}}$
from which $\Phi_{\mathrm{dyn}}$ is compiled.

\begin{figure}[t]
  \centering
  \resizebox{\linewidth}{!}{%
  \begin{tikzpicture}[auto, >=stealth, thick,
    block/.style={draw=blue!60, fill=blue!5, rectangle, rounded corners, minimum width=3.1cm, minimum height=0.95cm, align=center, thick},
    model/.style={draw=orange!60, fill=orange!5, rectangle, rounded corners, minimum width=3.1cm, minimum height=0.95cm, align=center, thick},
    arrow/.style={->, >=stealth, thick, draw=black!70},
    dashedarrow/.style={->, >=stealth, dashed, thick, draw=blue!80}
  ]
    \begin{scope}
      \node (iorA) [block] at (0,2.4) {Input/output\\requirement $R$};
      \node (fsrA) [block] at (0,0.8) {Functionality cotyledon\\$\mathrm{FSR}(R)$};
      \node (drefA) [model] at (0,-0.8) {Chosen design\\$D_{\mathrm{ref}}$};
      \node (implA) [model] at (4.0,-0.8) {Buildable resultant\\$Z_{\mathrm{impl}}$};
      \path[arrow] (iorA) -- node[right, font=\footnotesize] {generates} (fsrA);
      \path[arrow] (fsrA) -- node[right, font=\footnotesize] {choose} (drefA);
      \path[arrow] (implA) -- node[above, font=\footnotesize]
        {$h:Z_{\mathrm{impl}}\to Z_{\mathrm{ref}}$} (drefA);
      \node[inner sep=7pt, fit=(iorA)(fsrA)(drefA)(implA)] (panelAbox) {};
      \node[below=0.2cm of panelAbox, font=\headingfont\bfseries] {A. Classical Constructive T3SD};
    \end{scope}
    \begin{scope}[xshift=10.0cm]
      \node (iorB) [block] at (0,2.4) {Input/output\\requirement $R$};
      \node (phiR) [block] at (0,0.8) {Boundary fragment\\$\Phi_R$};
      \node (afsr) [block] at (0,-0.8) {Assertional space\\$\mathrm{AFSR}(\Phi_R)\simeq\mathrm{FSR}(R)$};
      \node (drefB) [model] at (4.0,2.4) {$D_{\mathrm{ref}}$ with system\\$Z_{\mathrm{ref}}$};
      \node (phiDyn) [block] at (4.0,0.8) {Dynamics fragment\\$\Phi_{\mathrm{dyn}}(Z_{\mathrm{ref}})$};
      \node (implB) [model] at (4.0,-0.8) {Buildable resultant\\$Z_{\mathrm{impl}}$};
      \path[dashedarrow] (iorB) -- node[right, font=\footnotesize] {compile} (phiR);
      \path[dashedarrow] (phiR) -- node[right, font=\footnotesize] {generates} (afsr);
      \path[arrow] (drefB) -- node[right, font=\footnotesize] {compile} (phiDyn);
      \path[dashedarrow] (implB) -- node[right, font=\footnotesize] {check} (phiDyn);
      \node[inner sep=8pt, fit=(iorB)(phiR)(afsr)(drefB)(phiDyn)(implB)] (panelBbox) {};
      \node[below=0.2cm of panelBbox, font=\headingfont\bfseries] {B. Assertional presentation};
    \end{scope}
  \end{tikzpicture}%
  }
  \caption{The IOR generates functionality in both presentations. A chosen functional
  design supplies $Z_{\mathrm{ref}}$ for the separate implementation check.}
  \label{fig:comparison}
\end{figure}

\subsection{Characteristic Boundary Fragment}
\label{sec:trace-observability}
Let $R=(OLR,IR,ITR,OR,OTR,ER)$ be any Wymore IOR, and let $T_R$ be the time scale selected by $OLR$. Its characteristic boundary fragment is
\[
  \Phi_R=(T_R,ITR,A_R),
  \qquad
  A_R(f,g)\;\Longleftrightarrow\;g\in ER(f).
\]
This fragment captures complete input and output trajectories. It retains finite or
infinite operational length, any admissible-input set, and the full
eligible-output relation, including nondeterministic relations.

For $Z=(S,I,O,\delta,\rho)$, initial state $s_0$, and nonempty $T\subseteq T_R$, define $Z,s_0,T\models_{IO}\Phi_R$ when, for every $f\in ITR$ and every complete input $h$ agreeing with $f$ on $T_R$, there exists $g$ with $A_R(f,g)$ such that the output generated by $Z$ on $h$ agrees with $g$ on $T$. Let $\mathrm{AFSR}(\Phi_R)$ package $Z$, $s_0$, $T$, and a proof of this assertional satisfaction relation.

Temporal assertions are useful only if they preserve the functional designs
admitted by the original IOR:

\begin{theorem}[Characteristic boundary fragment]
\label{thm:boundary-characteristic}
For every Wymorian IOR $R$,
\[
 Z,s_0,T\models_{IO}\Phi_R
 \quad\Longleftrightarrow\quad
 SatisfiesIOR(Z,R,s_0,T).
\]
Consequently there is a canonical bijection
\[
  \mathrm{AFSR}(\Phi_R)\simeq\mathrm{FSR}(R)
\]
that fixes the system, designated state, and time scale.
\end{theorem}

\begin{proof}[Proof sketch]
Substitute $T_R$, $ITR$, and $A_R(f,g)\Leftrightarrow g\in ER(f)$ into the
assertional satisfaction definition. Its clauses are those of
$SatisfiesIOR$. The bijection preserves the design data and converts the
satisfaction certificate in either direction.
\end{proof}

The fragment omits a separate $OTR$ field because Wymore's IOR axioms define
that set through the ranges of $ER$. The theorem covers finite and infinite
horizons, restricted input sets, and deterministic or relational eligible
outputs. Finite ordinary-LTL formulas are available when the trajectory
relation has a suitable presentation; Section~\ref{sec:case-basics} gives
three examples.

Implementation checking uses a second construction. Choose
$D_{\mathrm{ref}}\in\mathrm{FSR}(R)$ and write
$Z_{\mathrm{ref}}:=D_{\mathrm{ref}}.Z$. The transition and readout tables of
$Z_{\mathrm{ref}}$ generate $\Phi_{\mathrm{dyn}}$. Thus $\Phi_R$
characterizes functional designs for $R$, while $\Phi_{\mathrm{dyn}}$
characterizes systems that realize the selected functional design.

Recovering the maps from the fragment at every state requires runs starting from every state
of $Z_{\mathrm{impl}}$. Restricting the initial-state class leaves unreachable
states unconstrained. For a designated initial state, the result applies to the
subsystem it generates. The statements below use all states unless a reachable
subsystem is named.

\begingroup
\renewcommand{\arraystretch}{1.5}
\begin{tabularx}{\linewidth}{@{}P{0.28\linewidth}YY@{}}
    \rowcolor{tableheader}
    \headingfont\bfseries Feature & \headingfont\bfseries Constructive & \headingfont\bfseries Assertional \\
    \addlinespace[4pt]
    Functionality & $R$ generates $\mathrm{FSR}(R)$ & $R$ compiles to $\Phi_R$ and $\mathrm{AFSR}(\Phi_R)\simeq\mathrm{FSR}(R)$ \\
    \addlinespace[4pt]
    Implementation evidence & Explicit homomorphism $h$ & Satisfaction of compiled $\Phi_{\mathrm{dyn}}$ \\
    \addlinespace[4pt]
    Guarantee & Trajectory simulation under $h$ & The same simulation, by the bi-implication \\
    \addlinespace[4pt]
    Maintained artifact & Pair-specific maps $h_S,h_I,h_O$ & Reference-indexed fragment $\Phi_{\mathrm{dyn}}$ \\
\end{tabularx}
\captionof{table}{The proposed workflow changes the maintained evidence, not Wymore's functionality or implementation semantics.}
\label{tab:comparison}
\endgroup
\vspace{6pt}


\subsection{Dynamics-Encoding Fragment}
\label{sec:dynamics-encoding-fragment}

For the classical implementation question, the relevant reference is
$Z_{\mathrm{ref}}$, the system component of $D_{\mathrm{ref}}$. We compile its
state-transition and readout tables into a \emph{dynamics-encoding} fragment
$\Phi_{\mathrm{dyn}}$. Throughout, ``homomorphism'' means the surjective,
transition- and readout-preserving triple $h=(h_S,h_I,h_O)$ of
\citet{Wymore1993}; the appendix states the indexed definition.

Table~\ref{tab:phi-dyn-fragment} presents the fragment that makes the
bi-implication hold. Clauses are compiled from $Z_{\mathrm{ref}}$ and checked
on $Z_{\mathrm{impl}}$ under surjective maps $h_S,h_I,h_O$, or under identity
when labels are shared. At each tick, $\mathit{refState}(s)$ means that
$h_S$ maps the implementation state to reference state $s$. Inputs and outputs
are interpreted similarly through $h_I$ and $h_O$.

\begingroup
\renewcommand{\arraystretch}{1.35}
\footnotesize
\begin{tabularx}{\linewidth}{@{}A{0.22\linewidth} A{0.78\linewidth}@{}}
    \rowcolor{tableheader}
    \headingfont\bfseries Element & \headingfont\bfseries Dynamics-encoding fragment $\Phi_{\mathrm{dyn}}$ \\
    \addlinespace[2pt]
    Operators & $G$ (always), $X$ (next), $\rightarrow$, $\wedge$; \textbf{not} $F$ (eventually) \\
    \addlinespace[2pt]
    Atoms & $\mathit{state}(s)$ / $\mathit{refState}(s)$; $\mathit{input}(i)$, $\mathit{input}(\mathrm{none})$; $\mathit{output}(o)$, $\mathit{output}(\mathrm{none})$ \\
    \addlinespace[2pt]
    Open readout & $G\bigl(\mathit{refState}(s) \rightarrow \mathit{output}(o)\bigr)$ whenever $\rho_{\text{ref}}(s)=o$ \\
    \addlinespace[2pt]
    Closed readout & $G\bigl(\mathit{refState}(s) \rightarrow \mathit{output}(\mathrm{none})\bigr)$ whenever $\rho_{\text{ref}}(s)=\mathrm{none}$ \\
    \addlinespace[2pt]
    Autonomous step & $G\bigl((\mathit{refState}(s)\wedge\mathit{input}(\mathrm{none})) \rightarrow X(\mathit{refState}(\delta_{\text{ref}}(s,\mathrm{none})))\bigr)$ \\
    \addlinespace[2pt]
    Guided step & $G\bigl((\mathit{refState}(s)\wedge\mathit{input}(i)) \rightarrow X(\mathit{refState}(\delta_{\text{ref}}(s,i)))\bigr)$ for each $i\in I_{\text{ref}}$ \\
    \addlinespace[2pt]
    Compiled from & System component $Z_{\mathrm{ref}}$ of $D_{\mathrm{ref}}$ \\
    \addlinespace[2pt]
    Checked on & Buildable system $Z_{\text{impl}}$ under surjective $h_S,h_I,h_O$ (or identity on shared labels) \\
    % \addlinespace[2pt]
    % Outside fragment & $F$; incomplete $\Phi_{\mathrm{dyn}}$; output-table-only; arbitrary FO-LTL outside the clause schemas \\
\end{tabularx}
\captionof{table}{Dynamics-encoding fragment $\Phi_{\mathrm{dyn}}$ for which assertional satisfaction is equivalent to existence of a Wymore system homomorphism.}
\label{tab:phi-dyn-fragment}
\endgroup
\vspace{4pt}

LTL inspires the notation for this fragment; however, predicate-indexed
clauses allow infinite or real-valued $S$, $I$, and $O$. Closed readouts and
autonomous steps remain are explicitly allowed; finite-state instances can be flattened for automated checking.

A property that we would like is that an assertional check must provide the same implementation evidence as the Wymore homomorphism it replaces:

\begin{theorem}[Dynamics-fragment characterization]
\label{thm:dynamics-characterization}
Fix $D_{\mathrm{ref}}\in\mathrm{FSR}(R)$ and its system
$Z_{\mathrm{ref}}$. Then
\[
  Z_{\text{impl}} \models_{\mathrm{dyn}} \Phi_{Z_{\text{ref}}}
  \;\;\Longleftrightarrow\;\;
  \exists\, h : Z_{\text{impl}} \to Z_{\text{ref}}\ \text{Wymore system homomorphism}.
\]
\end{theorem}

\begin{proof}[Proof sketch]
A Wymore homomorphism satisfies the readout and transition equations at each
implementation state. Applying those equations along a run yields the four
clause families of Table~\ref{tab:phi-dyn-fragment}. Conversely, start a run
from an arbitrary implementation state and inspect the clauses at tick zero.
They recover the readout equation and each autonomous or guided transition
equation. Appendix~\ref{app:hom-biimp} gives the indexed argument.
\end{proof}

A conforming candidate must, of course, also be buildable before it supports a T3SD
implementation decision:

\begin{corollary}[Cotyledon implementability]
\label{cor:cotyledon-implementability}
If $Z_{\mathrm{impl}}\models_{\mathrm{dyn}}\Phi_{Z_{\mathrm{ref}}}$ and
$Z_{\mathrm{impl}}$ is the resultant of a buildable design in $TYR$, then
$D_{\mathrm{ref}}$ and that buildable design form an element of
$\mathrm{ISR}(R,TYR)$.
\end{corollary}

\begin{proof}[Proof sketch]
Theorem~\ref{thm:dynamics-characterization} gives the homomorphism, hence a witness of
implementation. Pair this with the given functional and
buildable designs. Requirement inheritance uses the homomorphism for identity
boundary maps and transport for general maps.
\end{proof}

Shared design vocabularies should turn failed assertions into direct table
discrepancies:

\begin{proposition}[Shared-table specialization]
\label{prop:shared-table}
If $Z_{\mathrm{ref}}$ and $Z_{\mathrm{impl}}$ share $(S,I,O)$ and use
the identity state, input, and output maps, satisfaction of
$\Phi_{Z_{\mathrm{ref}}}$ is
equivalent to pointwise agreement of their transition and readout functions.
\end{proposition}

\begin{proof}[Proof sketch]
With canonical labels the three projection maps are identities. The readout
clauses give equality of $\rho$ at each state, and the transition clauses give
equality of $\delta$ for every optional input. The converse follows by
substitution in the clauses.
\end{proof}

Cross-type elaborations use the hom-relative reading of
Theorem~\ref{thm:dynamics-characterization}. Note that finite propositional projections can support automated checking when the relevant spaces are finite.

It is worth noting that $\Phi_{\mathrm{dyn}}$ is a well-known fragment of temporal logic. Each clause is
$G$ of a Horn implication whose consequent lies under at most one $X$, making
satisfaction a per-state safety condition. Compiling a transition table into a
formula is a characteristic-formula construction, and the maps required by a
Wymore homomorphism are bounded morphisms. This identification connects T3SD
implementation checking to established theory and tools
(Section~\ref{sec:case-solver}).

\subsection{Fragment Invariance and Reuse}
\label{sec:invariance}

The compiled fragment is determined by $Z_{\mathrm{ref}}$. A witnessing
homomorphism also depends on the implementation and on both encodings. This
difference determines what can be reused when a design changes.

Representational changes over a system lifecycle should not invalidate an
established behavioral result.

\begin{theorem}[Encoding invariance]
\label{thm:encoding-invariance}
Replacing either $Z_{\mathrm{ref}}$ or $Z_{\mathrm{impl}}$ by an isomorphic
system preserves satisfaction of the compiled dynamics fragment. The same
holds for T3SD copies that rename or permute ports.
\end{theorem}

\begin{proof}[Proof sketch]
Compose the witnessing homomorphism with the isomorphism in the required
direction. The inverse isomorphism proves the reverse implication. T3SD copies
carry the corresponding bijections on state and boundary indices, so the same
composition argument applies.
\end{proof}

Component-level evidence must lift to assembled resultants if verification is
to scale with system architecture.

\begin{theorem}[Elaboration and coupling]
\label{thm:elaboration-coupling}
Dynamics-fragment conformance composes through successive homomorphic
realisations on either side. For two coupling recipes with the same free
boundary, port-preserving component homomorphisms induce a homomorphism between
their resultants.
\end{theorem}

\begin{proof}[Proof sketch]
Composition of the state, input, and output maps preserves surjectivity and the
homomorphism equations. For coupled systems, take the componentwise product of
the state maps. Port preservation makes the internal-wire values agree, so the
product map preserves the resultant transition and readout functions. Standard
T3SD coupling results supply the corresponding resultant construction
\citep{Wymore1993}.
\end{proof}

Reindexing components, flattening nested recipes, and removing
boundary-irrelevant components are structural instances of
Theorem~\ref{thm:elaboration-coupling}. Section~\ref{sec:case-machine-lift}
illustrates component elaboration by changing the instruction table's phase
encoding.

Together, these results replace one monolithic obligation with a hierarchy of checks.
Each component can be checked against its reference fragment; the coupling
theorem lifts the component results when the required coupling witness is
supplied; subsequent elaborations reuse the established fragment. Finite
components admit propositional checking, while infinite components may use
symbolic proofs or finite abstractions. The RLE case proves conformance of the
full resultant directly and illustrates fragment reuse on its instruction
table. The underlying search problem remains difficult in general, while
decomposition and reuse provide a route to larger architectures.

A characteristic fragment should distinguish finite behavior, not merely admit
coarse mutual projections.

\begin{theorem}[Finite canonicity]
\label{thm:finite-canonicity}
Let $Z_1$ and $Z_2$ be finite systems. If each satisfies the dynamics fragment
compiled from the other, then $Z_1$ and $Z_2$ are isomorphic.
\end{theorem}

\begin{proof}[Proof sketch]
Mutual satisfaction yields surjective homomorphisms in both directions by
Theorem~\ref{thm:dynamics-characterization}. Finiteness forces the state,
input, and output maps to be bijective. Their inverses preserve the same
transition and readout equations, giving an isomorphism.
\end{proof}

Theorems~\ref{thm:encoding-invariance}--\ref{thm:finite-canonicity} support
incremental verification while preserving Wymore's implementation semantics.
Composition lifts established component results, and finite presentations turn
local conformance obligations into generated decision problems.


\subsection{Property Checking}
\label{sec:automated-checking}

Property checking of $Z_{\text{impl}} \models_{\mathcal{L}} \Phi$ may use
finite-horizon unrolling, finite-state abstraction, or theory-specific solvers
on bounded queries \citep{BaierKatoen2008,deMouraZ32008}. Unbounded FO-LTL
over infinite or real-valued coordinates is undecidable in general. Finite
projections, bounded horizons, restricted input classes, and symbolic proofs
cover different parts of the resulting verification problem.

We leave mechanical verification to future work, considering only the theoretical aspects here.
